% ============================================================================
% CHAPTER 7: REGIONAL SERVICE DELIVERY PLAN
% ============================================================================

\chapter{Regional Service Delivery Plan}

\section{Regional Strategy Overview}

The expansion plan is straightforward: start where the revenue is strongest, prove the model, then push outward. Dhaka comes first, Chittagong follows in Month 3, and Khulna and Sylhet open in Years 2 and 3 as revenue supports the investment.

\section{Region 1: Dhaka \& Central (Priority 1 - Year 1)}

\subsection{Geographic Coverage}

Dhaka and its satellite cities --- Narayanganj, Gazipur, Tangail --- form the country's primary economic engine, with the highest concentration of government agencies, real estate developers, and infrastructure projects.

\subsection{Market Characteristics}

\begin{itemize}
    \item \textbf{Key Clients:}
    \begin{itemize}
        \item Dhaka City Corporation
        \item Major real estate developers (BASHUNDHARA, Nasha, Concord)
        \item Government ministries (Land, Public Works, Disaster Management)
        \item International development organizations (ADB, World Bank projects)
    \end{itemize}

    \item \textbf{Service Demand:}
    \begin{itemize}
        \item Urban planning \& infrastructure mapping
        \item Real estate site selection
        \item Land survey \& documentation
        \item Building footprint extraction
        \item Infrastructure project design support
    \end{itemize}

    \item \textbf{Market Size:} BDT 40-50 lakhs annual opportunity
\end{itemize}

\subsection{Implementation Timeline}

\begin{itemize}
    \item \textbf{Month 1:} Office setup, team hiring, equipment procurement
    \item \textbf{Month 2:} Regulatory approvals, pilot projects
    \item \textbf{Month 3:} Government project targeting begins
    \item \textbf{Month 6:} Real estate client acquisition starts
    \item \textbf{Month 12:} Establish operational excellence, expand to Region 2
\end{itemize}

\subsection{Success Metrics}

\begin{itemize}
    \item 5-8 projects completed
    \item 2-3 case studies for portfolio
    \item 5-10 government/corporate client relationships
    \item BDT 40-50 lakhs revenue
\end{itemize}

\section{Region 2: Chittagong \& Southeast (Priority 2 - Year 1)}

\subsection{Geographic Coverage}

Chittagong and the southeastern belt --- Cox's Bazar, Rangamati, Bandarban --- is Bangladesh's second economic centre, with terrain that makes LiDAR considerably more valuable than photogrammetry alone.

\subsection{Market Characteristics}

\begin{itemize}
    \item \textbf{Key Clients:}
    \begin{itemize}
        \item Chittagong Port Authority
        \item Tourism development agencies
        \item Hill tract development corporations
        \item Coastal/beach conservation projects
        \item Disaster management for cyclone preparedness
    \end{itemize}

    \item \textbf{Service Demand:}
    \begin{itemize}
        \item Port infrastructure mapping
        \item Coastal zone management
        \item Tourism site planning
        \item Forest inventory (hill tracts)
        \item Landslide risk mapping
        \item Cyclone flood modeling
    \end{itemize}

    \item \textbf{Unique Capability:} Hilly terrain is where LiDAR earns its premium. Photogrammetry struggles with dense vegetation and complex elevation; LiDAR cuts through both.

    \item \textbf{Market Size:} BDT 20-25 lakhs annual opportunity
\end{itemize}

\subsection{Implementation Timeline}

\begin{itemize}
    \item \textbf{Month 3:} Regional hub opening in Chittagong
    \item \textbf{Month 4:} Port Authority project targeting
    \item \textbf{Month 6:} Hilltract survey pilot projects
    \item \textbf{Month 12:} Established regional presence
\end{itemize}

\subsection{Success Metrics}

\begin{itemize}
    \item 5-8 projects (mix of government and private)
    \item Specialized hill-terrain expertise demonstrated
    \item BDT 20-25 lakhs revenue
    \item Regional hub established with 3-4 staff
\end{itemize}

\section{Region 3: Khulna \& Southwest (Priority 3 - Year 2)}

\subsection{Geographic Coverage}

The southwest delta --- Khulna, Barisal, Satkhira, Jessore --- is agricultural heartland and home to the NGO networks that make the farmer application platform viable.

\subsection{Market Characteristics}

\begin{itemize}
    \item \textbf{Key Clients:}
    \begin{itemize}
        \item Agricultural cooperatives \& farming groups
        \item Shrimp farming operations
        \item Sundarbans management authority
        \item NGOs in development and conservation
        \item Environmental research institutions
    \end{itemize}

    \item \textbf{Service Demand:}
    \begin{itemize}
        \item Precision agriculture mapping
        \item Crop health monitoring (NDVI)
        \item Land use classification
        \item Water resource mapping
        \item Soil health assessment
        \item Mangrove/Sundarbans monitoring
    \end{itemize}

    \item \textbf{Unique Advantage:} BRAC, ASA, and Practical Action operate heavily here. Those existing networks become the distribution channel for the farmer apps.

    \item \textbf{Market Size:} BDT 30-35 lakhs annual opportunity
\end{itemize}

\subsection{Implementation Timeline}

\begin{itemize}
    \item \textbf{Month 6, Y2:} Regional hub opening in Khulna
    \item \textbf{Month 9, Y2:} NGO partnership pilots begin
    \item \textbf{Month 12, Y2:} First 5K farmers on Data Science platform
    \item \textbf{Month 18, Y2:} Expanded partnership with 2-3 NGOs
\end{itemize}

\subsection{Success Metrics}

\begin{itemize}
    \item 10-15 projects
    \item 5K+ farmers reached through NGO channels
    \item 2-3 NGO strategic partnerships
    \item BDT 30-35 lakhs revenue
    \item Regional hub with 3-4 staff
\end{itemize}

\section{Region 4: Sylhet \& Northeast (Priority 4 - Year 3)}

\subsection{Geographic Coverage}

Sylhet, Moulvibazar, Sunamganj --- tea country. The focus here is research partnerships, speciality crop monitoring, and biodiversity work.

\subsection{Market Characteristics}

\begin{itemize}
    \item \textbf{Key Clients:}
    \begin{itemize}
        \item Tea garden operators (80+ major estates)
        \item Forest research institutions
        \item Climate adaptation projects
        \item Agricultural universities
        \item Environmental research centers
    \end{itemize}

    \item \textbf{Service Demand:}
    \begin{itemize}
        \item Tea plantation health monitoring
        \item Forest inventory \& conservation
        \item Climate vulnerability mapping
        \item Agricultural research data
        \item Ecosystem monitoring
        \item Species mapping
    \end{itemize}

    \item \textbf{Unique Focus:} Research partnerships; specialized crops (tea); biodiversity

    \item \textbf{Market Size:} BDT 15-20 lakhs annual opportunity
\end{itemize}

\subsection{Implementation Timeline}

\begin{itemize}
    \item \textbf{Month 6, Y3:} Regional hub opening (partnered model)
    \item \textbf{Month 9, Y3:} Tea estate survey projects
    \item \textbf{Month 12, Y3:} Research partnerships established
\end{itemize}

\subsection{Success Metrics}

\begin{itemize}
    \item 5-8 projects
    \item Tea plantation monitoring platform established
    \item 2-3 research partnerships
    \item BDT 15-20 lakhs revenue
    \item Regional presence with 2-3 staff
\end{itemize}

\section{Regional Hub Expansion Timeline}

\begin{table}[h!]
    \centering
    \caption{Regional Hub Development Schedule}
    \label{tab:hub_expansion}
    \begin{tabularx}{\textwidth}{l | X | X | X | X}
        \hline
        \theader{Region} & \theader{Launch} & \theader{Staff} & \theader{Investment} & \theader{Revenue Target} \\
        \hline
        Dhaka (HQ) & M1, Y1 & 8-10 & 50L & 40-50L \\
        \hline
        Chittagong & M3, Y1 & 3-4 & 15L & 20-25L \\
        \hline
        Khulna & M6, Y2 & 3-4 & 15L & 30-35L \\
        \hline
        Sylhet & M6, Y3 & 2-3 & 10L & 15-20L \\
        \hline
        \textbf{TOTAL} & & \textbf{16-21} & \textbf{90L} & \textbf{105-130L} \\
        \hline
    \end{tabularx}
\end{table}

\section{Inter-Regional Synergies}

\subsection{Knowledge Sharing}

\begin{itemize}
    \item Consistent operating procedures across all four hubs mean a project team from Dhaka can plug into a Chittagong job without a learning curve
    \item Learnings from each project type flow back to all hubs through monthly team reviews
    \item All heavy data processing runs through Dhaka, keeping infrastructure costs consolidated
\end{itemize}

\subsection{Equipment Sharing}

\begin{itemize}
    \item Drone units travel --- a busy period in Dhaka doesn't leave Chittagong grounded
    \item Processing workloads are balanced across hubs based on capacity
    \item Backup equipment is shared across the network for redundancy
\end{itemize}

\subsection{Project Allocation}

\begin{itemize}
    \item Large national projects coordinated across hubs
    \item Regional specialists support neighboring regions
    \item Emergency response (cyclones, disasters) coordinated nationally
\end{itemize}

% ============================================================================
% END OF CHAPTER 7
% ============================================================================
